% ==============================================================================
\section{Introduction}
% ==============================================================================

Global food waste represents an urgent sustainability and resource-management challenge. The United Nations Environment Programme (UNEP) \textit{Food Waste Index Report 2024} estimates that 1.05 billion tonnes of food were wasted across retail, food service, and household sectors in 2022, with households accounting for roughly 60\% of the total \cite{ref1}. In everyday household cooking, users must track what ingredients remain across refrigerators, freezers, and pantries, how much quantity is left, and which items require urgent consumption. Without synchronized inventory visibility and actionable meal planning, households face recurring over-purchasing and avoidable spoilage.

Although large language models (LLMs) and conversational recipe assistants are widely used, general-purpose chatbots operate without access to a household's live inventory ledger. Consequently, they often suggest dishes requiring missing ingredients, overlook near-expiry stock, or fail to justify recommendations with traceable evidence. Moreover, food date labels such as ``Use By'' and ``Best Before'' are frequently misunderstood by consumers, causing either premature disposal of safe food or unsafe consumption of perishable items \cite{ref2}, \cite{ref8}, \cite{ref9}.

In this paper, we formulate household meal planning as generating verifiable, serving-scaled meal recommendations under known inventory, date-label constraints, and potentially incomplete storage metadata. For example, if a recipe for two servings of tomato scrambled eggs requires 4 eggs and 2 tomatoes, while SQLite inventory holds 3 eggs and 2 tomatoes, a reliable assistant must compute and report the exact deficit (1 missing egg) rather than hallucinating sufficiency, and must distinguish immediately cookable dishes from dishes prioritizing near-expiry items.

Existing solutions typically isolate these capabilities: inventory tables report ``what is on hand,'' reminder apps flag ``what is nearing a date label,'' and recipe tools describe ``how to cook a dish.'' Bridging them into a single household decision requires verifying quantities across compatible units, applying conservative safety exclusions prior to recommendation, and logging post-cooking outcomes without conflating recommendation with actual consumption. While recent prototypes such as MealMind explore computer-vision ingredient capture and LLM meal planning \cite{ref3}, there remains a need for lightweight, reproducible architectures that ground manual inventories in deterministic arithmetic and auditable tool traces.

To address these challenges, we propose \textit{Food Guardian}, a hybrid AI agent scoped for single-household deployment. Its main contributions are:
\begin{enumerate}
    \item \textbf{Hybrid Tool-Calling Architecture:} We decouple natural language intent parsing and orchestration (Gemini via LangGraph) from deterministic computation (Python services over SQLite), ensuring reproducible serving scaling, unit normalization, shortage calculation, and safety pre-screening.
    \item \textbf{Dual-Mode Inventory-Aware Ranking:} We implement \texttt{ready\_now} and \texttt{expiry\_first} modes (retaining \texttt{balanced} compatibility) to prioritize either zero-shopping feasibility or near-expiry utilization with explicit deficit reporting.
    \item \textbf{Conservative Safety Screening and Source-Attributed Knowledge:} We enforce rule-based batch filtering (\texttt{eligible}, \texttt{review}, \texttt{blocked}) over SQLite records and integrate structured JSON plus ChromaDB RAG that exposes source verification status rather than fabricating safety assurances.
    \item \textbf{Transactional Analytics and OpenSpec Verification:} We separate read-only recommendation from user-confirmed consumption/waste logging in atomic SQLite transactions, and formalize requirements \texttt{FR-01}--\texttt{FR-10} via OpenSpec for auditable acceptance testing.
\end{enumerate}

The remainder of this paper is organized as follows. Section~\ref{sec:related_work} reviews AI food recommendation and household food waste studies. Section~\ref{sec:proposed_design} presents our design principles, hybrid architecture, data model, and core algorithms. Section~\ref{sec:experiment} details prototype implementation, OpenSpec verification, preliminary results, and the comparative evaluation plan. Section~\ref{sec:conclusion} concludes the paper.
